\section{Instruction Level Parallelism (ILP)}

\subsection{Pipelining}

\begin{definition}
\textbf{Pipelining} is a microarchitecture implementation technique where \textbf{multiple instructions are overlapped} in execution. It functions like an assembly line, dividing the processing of an instruction into a sequence of stages so that a different instruction can be processed in each stage simultaneously.
\end{definition}

\begin{important}
The core objective of pipelining is to increase \textbf{instruction throughput} (the number of instructions completed per unit of time). While it significantly improves throughput, it generally does not decrease the \textbf{latency} of a single instruction; in fact, the latency may increase slightly due to the overhead of pipeline registers.
\end{important}

\begin{theorem}
In an ideal scenario where all stages have equal delay and there are no dependencies, a pipeline with $k$ stages provides a $k$-fold speedup. However, in reality, the speedup is limited by:
\begin{itemize}
    \item \textbf{Unbalanced stages:} The clock period is determined by the slowest stage.
    \item \textbf{Sequencing overhead:} Pipeline registers add setup time and clock-to-Q delays.
    \item \textbf{Hazards:} Dependencies between instructions that prevent seamless overlapping.
\end{itemize}
\end{theorem}



\subsubsection{The 5-Stage Instruction Pipeline}

\begin{definition}
A standard RISC pipeline (like MIPS or LC-3) typically consists of five stages:
\begin{enumerate}
    \item \textbf{Instruction Fetch (IF):} Fetch the instruction from memory and update the PC.
    \item \textbf{Instruction Decode (ID):} Decode the instruction and read operands from the register file.
    \item \textbf{Execute (EX):} Perform an arithmetic/logic operation or calculate a memory address.
    \item \textbf{Memory Access (MEM):} Read from or write to the data memory.
    \item \textbf{Write-Back (WB):} Write the result back into the register file.
\end{enumerate}
\end{definition}

\begin{important}
To maintain the state of different instructions in different stages, \textbf{Pipeline Registers} (IF/ID, ID/EX, EX/MEM, MEM/WB) are placed between the stages. These registers carry the data and control signals forward through the pipeline every clock cycle.
\end{important}

\subsubsection{Pipeline Hazards and Dependencies}

\begin{definition}
A hazard is a situation that prevents the next instruction in the stream from executing in its designated clock cycle.
\begin{itemize}
    \item \textbf{Structural Hazards:} Conflicts for hardware resources (e.g., using a single memory for both instructions and data).
    \item \textbf{Data Hazards:} An instruction depends on the result of a previous instruction that has not yet completed.
    \item \textbf{Control Hazards:} The next address to fetch is unknown because a branch or jump has not been resolved.
\end{itemize}
\end{definition}

\begin{definition}
Data dependencies are primarily of three types:
\begin{itemize}
    \item \textbf{Flow Dependence (RAW):} Instruction $j$ needs data produced by instruction $i$.
    \item \textbf{Anti-Dependence (WAR):} Instruction $j$ writes to a location that instruction $i$ reads.
    \item \textbf{Output Dependence (WAW):} Both instructions write to the same location.
\end{itemize}
In a standard 5-stage in-order pipeline, only RAW hazards require hardware intervention.
\end{definition}

\subsubsection{Resolving Data Hazards}

\begin{lemma}
\textbf{Forwarding (Bypassing)} resolves many data hazards by routing the result of an instruction directly from the EX or MEM stage to the ALU input of a following instruction, bypassing the register file.
\end{lemma}

\begin{important}
Forwarding cannot resolve a \textbf{Load-Use Hazard} (where a load is immediately followed by an instruction using the loaded data). This requires a \textbf{Stall}. The hardware must:
\begin{enumerate}
    \item Disable the PC and IF/ID register updates (holding instructions in place).
    \item Flush the ID/EX register (inserting a "bubble" or NOP).
\end{enumerate}
\end{important}

\subsubsection{Control Hazards and Branching}

\begin{definition}
\textbf{Control hazards} arise from the delay in determining the branch outcome. If the pipeline fetches the wrong instructions, it must \textbf{flush} those instructions once the branch is resolved.
\end{definition}

\begin{remark}
Common strategies for handling branches include:
\begin{itemize}
    \item \textbf{Static Prediction:} For example, "Always Predict Not Taken."
    \item \textbf{Dynamic Prediction:} Using a Branch History Table to predict based on past execution.
    \item \textbf{Early Branch Resolution:} Moving the branch comparison to the Decode stage to reduce the misprediction penalty.
\end{itemize}
\end{remark}

\subsubsection{Advanced Pipeline Issues}

\begin{important}
\textbf{Precise Exceptions:} A pipelined processor must be able to stop execution on an exception-causing instruction while ensuring all prior instructions complete and no subsequent instructions have modified the architectural state. This is challenging because instructions finish at different times.
\end{important}

\begin{lemma}
In systems with multi-cycle functional units (e.g., a fast integer ALU and a slow divider), instructions can finish \textbf{out of order}. This "out-of-order completion" can break the sequential semantics of the ISA if an exception occurs in the middle of a stream.
\end{lemma}

\begin{definition}
\textbf{Fine-Grained Multithreading:} A technique where the processor switches to a different thread every clock cycle. This hides the latency of data and control hazards by ensuring that instructions currently in the pipeline are independent of one another.
\end{definition}


\subsubsection{Precise Exceptions and Interrupts}

In a pipelined Architecture, not all instructions take the same amount of time in the execution stage. 
So what happens if a previous instruction causes an exception and the next instruction has already been written back to the register file? 
We need to make sure that the state of the processor is consistent with the ISA. 

\begin{definition}
    Unplanned changes in program exection due to internal problems are called \textbf{exceptions} while due to external events they are called \textbf{interrupts}.
\end{definition}

\begin{important}
    When an exception/interrupt occurs all previous instructions should be retired whil no later instructions should be allowed to proceed nor be retired.
\end{important}

Hence when the oldest instruction ready to be retired is detected to have caused an exception we have to:
\begin{itemize}
    \item Ensure architectural state is consistent (REG-File, PC, Memory)
    \item Flush all younger instructions in the pipeline
    \item Save PC and REG
    \item Redirect to exception handler
\end{itemize}

\begin{remark}
    In Single-Cycle machines this is easy to handle as we only have one instruction in the pipeline at a time which ensures sequential execution semantics.
\end{remark}

\begin{remark}
    Note that we could make every operation take the same amount of time in the pipeline, this would ensure sequential execution semantics but would be very inefficient.
\end{remark}


\begin{remark}
    The idea behind a \textbf{reorder-buffer} is that we complete instructions out of order but reorder them before making the results visible to the architectural state.
\end{remark}

\begin{definition}
    A \textbf{Reorder-Buffer (ROB)} is a hardware structure that stores the results of instructions as they complete. It allows instructions to be retired in-order, even if they complete out-of-order.
\end{definition}

\begin{lemma}
    Reordering-Buffer uses a valid bit in the REG to indicate whether the result is valid, if not the REG stores a pointer to the ROB entry where the result will be stored once the instruction completes (ROB also uses a valid bit to indicate readiness).
\end{lemma}

This allows us to write out of order to the REG and ensure in order retirement where we can check in-order for any exceptions.

\subsection{Out of Order Execution}

\begin{remark}
    An instruction that uses operands of a previous instruction is called a \textbf{dependent instruction}. If such a preceding instructions result are not yet available, the dependent instruction must wait for the result to be computed, which stalls an \textbf{in-order pipeline}.
\end{remark}

\begin{definition}
    \textbf{Out-of-order execution (OoO)} is a technique that allows a processor to execute instructions in a different order than they appear in the program. This is done to improve performance.
\end{definition}

\begin{remark}
    OoO tolerates longer latencies by executing other independent instructions between dependent ones that would otherwise stall the pipeline. This improves instructions per cycle (IPC).
\end{remark}

\begin{definition}
    The \textbf{dataflow principle} states that the execution order of instructions should be determined by the availability of their operands rather than their program order.
\end{definition}

\begin{remark}
    Typically OoO is only done in the core. Hence instructions are \textbf{fetched and decoded in order} then they are \textbf{executed out of order} and the results are \textbf{written back in order}.
\end{remark}

\subsubsection{Mechanisms for OoO}

\begin{definition}
    \textbf{Register-Renaming} eliminates false dependencies (WAW and WAR), which occure because of the reuse of registernames, by decoupling the register names from the actual physical storage locations.
\end{definition}

\begin{definition}
    \textbf{Reservation Stations (RS)} are buffers that store instructions that are waiting for their operands to become available.
\end{definition}

\begin{definition}
    \textbf{Tag Broadcast} is a mechanism that allows the results of instructions to be broadcast to all reservation stations that are waiting for those operands by broadcasting identification tags of available results via a bus.
\end{definition}

\subsubsection{Tomasulo's Algorithm}

\begin{definition}
    \textbf{Tomasulo's Algorithm} is an algorithm for out-of-order execution that uses reservation stations and tag broadcasting to eliminate false dependencies and improve performance.
\end{definition}

\begin{algorithm}
decode(instruction)

if (free reservation station) {

    reservation station = instruction
    replace registers with tags from Register-Alias Table (RAT)

    if (operands ready) {

        execute to functional unit
        broadcast result and update RAT

    } else {
        observe common data bus for tags of operands
    }

} else {
    wait
}
\end{algorithm}

\begin{important}
    The instruction window size (how many instr. can be fetched/decoded and not yet retired) limits the scalability of Tomasulo's algorithm.
\end{important}

\begin{remark}
    Moder processors use the already discussed reorder buffer (ROB) to ensure in order commitment of results.
\end{remark}

\begin{remark}
    The usage of a \textbf{Scheduler} with reservation stations and a \textbf{reorder buffer} can be visualized as a camel with two humps.
\end{remark}

\begin{remark}
    OoO allows for larger latency tolerance and exploitation of Irregular parallelism but comes at the cost of increased complexity and power consumption.
\end{remark}

\subsubsection{Memory Dependencies}

\begin{example}
    Instruction $i$ writes to memory adress $x$ while the instruction $i+1$ reads from $x$. In out of order execution it can happen that $i+1$ finishes before $i$ which invalidates its result. Hence we need to handle memory dependencies.
\end{example}

\begin{definition}
    \textbf{Memory Disambiguation} is the process of determining whether a load/store is dependent on a previous load/store.
\end{definition}

\begin{definition}
    \textbf{Solutions to memory dependencies:}
    \begin{itemize}
        \item \textbf{Conservative: } Assume all load/stores are dependend and stall the pipeline to ensure correctness.
        \item \textbf{Agressive: } Assume no load/stores are dependent and pay the extra time/work when a false result is generated.
        \item \textbf{Intelligent: } Try and predict whether a load/store is dependent or not in order to minimize the penalty of false results.
    \end{itemize}
\end{definition}

\begin{remark}
    Observe that the intelligent approach is a trade-off between the conservative and aggressive approaches and ensures the best performance.
\end{remark}

\begin{definition}
    \textbf{Load/Store Queue (LSQ)} is a queue that stores load/store results that have not yet been committed to the reorder buffer.
\end{definition}

\begin{definition}
    \textbf{Store-to-Load Forwarding} is a mechanism that allows the result of a store instruction to be forwarded directly to a load instruction that is waiting for that result. This works by searching the LSQ for a matching address and directly copying the data from there.
\end{definition}


\subsection{Superscalar Execution}

\begin{definition}
    A $N$-\textbf{Superscalar Processor} is a processor that can issue more than one ($N$) instruction per clock cycle.
\end{definition}

\begin{remark}
    Note that superscalar execution is orthogonal to OoO execution. A processor can be in-order or out-of-order and scalar or superscalar.
\end{remark}

\begin{important}
    The main challenge of superscalar processors is to find independent instructions to execute in parallel. (For example trough Forwarding, Software-Optimizations, Speciulation or Multi-Threading)
\end{important}

\begin{remark}
    The \textbf{tradeoffs} of Superscalar execution are IPC vs Complexity.
\end{remark}

\subsection{Branch Prediction}

\begin{remark}
    \textbf{The Goal:} Branch prediction handles control dependencies. It attempts to answer the fundamental question: How do we exploit instruction-level parallelism in the presence of control flow?
\end{remark}

\begin{remark}
    \textbf{The Cost of Misprediction:} Branch mispredictions drastically reduce IPC in deep pipelines. Modern processors therefore target 95\%+ accuracy.
\end{remark}

\begin{definition}
    There are multiple different \textbf{Branch Types}: Conditional (if, while), Unconditional (goto), Indirect (register targets), \& Function calls/returns. \textit{Conditional} and \textit{indirect} are the hardest to predict.
\end{definition}

\begin{remark}
    Control dependencies are handled by:
    \begin{itemize}
        \item \textbf{Stalling} or \textbf{Predicting} the pipeline.
        \item \textbf{Delayed Branching:} Executing adjacent instructions unconditionally.
        \item \textbf{Multi-threading:} Executing other threads.
        \item \textbf{Predicated Execution:} Converting control to data-dependence.
        \item \textbf{Multi-path Execution:} Fetching all branches simultaneously.
    \end{itemize}
\end{remark}

\subsubsection{Static Branch Prediction}

\begin{definition}
    \textbf{Static Branch Prediction}: The branch outcome is predicted strictly at compile time. Because it is static, it cannot adapt to dynamic changes in branch behavior during runtime.
\end{definition}

\begin{remark}
    The simplest strategy is to always predict the same outcome (\textbf{Always Taken} or \textbf{Always Not-Taken}). This logic requires minimal hardware but yields low accuracy.
\end{remark}

\begin{definition}
    \textbf{BTFN (Backward Taken, Forward Not-Taken)}: Predicts backward branches as taken and forward branches as not-taken.
\end{definition}
\begin{remark}
    BTFN significantly improves prediction success by intelligently exploiting loop structures. Backward branches are typically loop iterations, which are statistically taken much more often than not.
\end{remark}

\begin{definition}
    \textbf{Profile-based Prediction}: Predicts the outcome based on branch behavior statistics gathered from a previous profiling run of the program.
\end{definition}

\begin{definition}
    \textbf{Program-based Prediction}: Uses compiler heuristics (e.g., assuming a check for a negative integer is an error path, hence not-taken).
\end{definition}

\begin{definition}
    \textbf{Programmer-based Prediction}: Relies on explicit pragmas or hints provided by the programmer directly in the source code.
\end{definition}

\begin{codeexample}
    // Pragmas as hints for the branch predictor
    if (likely(x)) {...}
    if (unlikely(error)) {...}
\end{codeexample}

\subsubsection{Dynamic Branch Prediction}

\begin{definition}
    \textbf{Dynamic Branch Prediction}: The branch outcome is predicted at runtime based on the execution history.
\end{definition}

\begin{definition}
    \textbf{Branch Target Buffer (BTB)}: A hardware cache that specifically stores the \textit{target addresses} of recently executed branches. When a branch instruction is fetched, the BTB is accessed to provide the target address immediately, avoiding pipeline stalls while waiting for the ALU to compute the target.
\end{definition}

\begin{definition}
    \textbf{Return Address Stack (RAS)}: Dedicated hardware for predicting function returns. It pushes the return address on a \texttt{call} instruction and pops it on a \texttt{ret}. It is highly accurate (>95\% with just an 8-entry stack) for predicting indirect returns.
\end{definition}

\begin{definition}
    \textbf{Indirect Branch Predictor}: Indirect branches map to multiple targets, making standard BTBs inaccurate. These are typically predicted using a dedicated history-based target predictor (e.g., indexing the BTB with the GHR XORed with the PC).
\end{definition}

\begin{definition}
    \textbf{Last-Time Predictor}: Predicts the outcome of a branch to be exactly the same as the outcome of the last time that specific branch was executed.
\end{definition}
\begin{remark}
    Last-Time predictors are simple (one bit per branch). They exploit long loops (e.g. 1000 iterations) with high accuracy but perform extremely poorly on short loops (e.g. 2 iterations) or on branches that alternate heavily.
\end{remark}

\begin{definition}
    \textbf{Two-Bit Counter (Bimodal Predictor)}: Prediction only changes after two consecutive mistakes. This introduces \textbf{hysteresis} (via cyclic Strong/Weak states), preventing the prediction from arbitrarily flipping due to a single anomalous outcome.
\end{definition}
\begin{remark}
    Two-Bit Counters significantly improve prediction accuracy compared to Last-Time predictors (yielding around 0.85-0.90 accuracy).
\end{remark}

\begin{definition}
    \textbf{Two-Level Prediction}: Attempts to correlate the outcome of a branch using history tracking:
    \begin{itemize}
        \item \textbf{Global History Register (GHR)}: Tracks the outcome of \textit{recent different branches} (global context).
        \item \textbf{Local History Register (LHR)}: Tracks the outcome of \textit{the same branch} in its previous iterations.
    \end{itemize}
\end{definition}

\begin{definition}
    \textbf{Gshare Predictor}: XORs the Global History Register (GHR) with the Program Counter (PC) to index the Pattern History Table (PHT), providing branch-specific context and drastically reducing interference.
\end{definition}

\begin{definition}
    \textbf{Agree Predictor} \& \textbf{Gskew Predictor}: Advanced techniques beyond Gshare to reduce interference. The \textbf{Agree Predictor} predicts whether the outcome will \textit{agree} or \textit{disagree} with a statically determined bias bit. The \textbf{Gskew Predictor} uses multiple PHTs with different hash indexing functions and resolves predictions via a majority vote.
\end{definition}

\begin{important}
    \textbf{Biased Branches \& Branch Filtering}: Filtering heavily biased branches (e.g., 99\% taken) to simpler predictors saves space in history tables for hard-to-predict branches.
\end{important}

\begin{important}
    \textbf{Tournament Predictor}: In many systems, branch prediction is not done using a single predictor. Instead, multiple different predictors run simultaneously, and a meta-predictor chooses whichever predictor has been the most accurate recently. 
\end{important}

\begin{remark}
    \textbf{Notable Historical \& Modern Implementations:}
    \begin{itemize}
        \item \textbf{Intel Pentium Pro}: First commercially successful out-of-order processor, utilizing two-level global prediction.
        \item \textbf{Alpha 21264}: Famously utilized a highly accurate Tournament Predictor.
        \item \textbf{Intel Pentium M}: Introduced specialized Loop and Jump predictors alongside traditional predictors.
        \item \textbf{AMD Zen 2 (and newer)}: Employs a highly accurate hybrid \textbf{Perceptron + TAGE} multi-level branch predictor.
    \end{itemize}
\end{remark}

\subsubsection{Advanced Predictors}

\begin{definition}
    \textbf{Perceptron Predictor}: Uses a single-layer neural network to predict the outcome of a branch.
\end{definition}

\begin{remark}
    Efficiently implemented as an \textit{adder tree} followed by a sign check. Achieves top-tier accuracy (0.95+) but \textbf{only learns linearly separable functions} (cannot learn XOR-type correlations).
\end{remark}

\begin{definition}
    \textbf{TAGE (TAgged GEometric history length)}: Uses multiple global history registers of varying, geometric lengths to predict the outcome of a branch.
\end{definition}

\begin{remark}
    TAGE predictors represent state-of-the-art branch prediction. They are more complex to implement than perceptron predictors but achieve the highest structural prediction accuracy available (0.97+).
\end{remark}

\subsubsection{Alternative Control Flow Handling}

\begin{definition}
    \textbf{Delayed Branching}: The $N$ instructions fetched after the branch (delay slots) execute unconditionally. \\
    \textit{Advantage}: Keeps pipeline full. \textit{Disadvantage}: Ties ISA to pipeline depth; hard to fill slots with useful independent instructions.
\end{definition}

\begin{definition}
    \textbf{Predicated Execution (e.g., CMOV)}: Eliminates branches via condition codes. An operation executes but only commits if the predicate is true.\\
    \textit{Advantage}: Eliminates mispredictions, allows larger basic blocks. \textit{Disadvantage}: Wastes energy fetching/executing paths that are discarded.
\end{definition}

